\section{从命题形状开始证明}
\label{sec:01-Mathematical-Language/03-Proof-Techniques/01}

\subsection{一张白纸面前的恐慌}

你被要求证明：两个奇数的和一定是偶数。

你手里只有一张白纸，一支笔，以及这句话本身。你试了几个数：\(3+5=8\)，偶数；\(7+11=18\)，偶数；\(13+21=34\)，还是偶数。你甚至用 Python 跑了前一万个奇数对，全对。可问题是：你不敢交卷。因为你知道——测试一万个例子，跟证明它对\textbf{所有}奇数对都成立，是两回事。

这就是证明的真正入口：不是``用什么技巧''，而是先搞清楚，眼前的命题到底长成什么逻辑形状。形状决定了你需要交付什么才算数。

\subsection{命题的七种形状，以及每种形状让你做什么}

当你说一句话被证明了，这句话本身就规定了什么叫``证明它''。逻辑结构不一样，证明的终点就不一样。下面七种是最常见的形状，每一种都对应着一种无法绕过的交付标准：

\begin{itemize}
\tightlist
\item
  \(\forall x\in D, P(x)\)——全称命题。你要拿一个完全任意的 \(x\)（只满足论域 \(D\) 的条件，不加任何额外限制），然后证明 \(P(x)\) 对它就成立。一旦你给 \(x\) 偷加了不该有的条件，你的证明就只涵盖了一部分对象，剩下的漏了。
\item
  \(P\to Q\)——蕴含。假设 \(P\) 成立，在这个暂借的前提下，推导出 \(Q\)。你不需要证明 \(P\) 本身是真的，只需要证明``如果 \(P\) 真，那 \(Q\) 一定真''。
\item
  \(P\land Q\)——合取。劈成两半，分别证 \(P\)、分别证 \(Q\)，两条线独立完成。
\item
  \(P\lor Q\)——析取。你至少得把其中一个分支证出来，或者分情况覆盖所有可能，证明不管走哪条路结论都成立。
\item
  \(A=B\)——集合相等。几乎所有集合相等的证明都走同一套路：先证 \(A\subseteq B\)，再证 \(B\subseteq A\)。两个方向，缺一不可。
\item
  \(\exists x, P(x)\)——存在命题。直接造一个具体的 \(x\) 出来，然后把 \(P(x)\) 验给它看。或者调用一个定理，让定理替你保证存在。
\item
  \(P\leftrightarrow Q\)——等价。两个方向都得走：\(P\to Q\) 和 \(Q\to P\)。少一个方向，你就只证明了半边。
\end{itemize}

这个清单的用处不在于背诵，而在于：当你面对一道证明题手足无措时，先把它翻译成上面某一种形状，你就立刻知道该往哪里使劲了。全称命题？抓起一个任意的对象。存在命题？开工构造。等价？准备走两遍。

\subsection{直接证明：让定义替你干活}

回到那两个奇数。现在我们知道命题的形状是 \(\forall a,b\in\text{奇数}, a+b\text{ 是偶数}\)——一个全称命题。那就任取两个奇数 \(a,b\)。接下来怎么办？

定义是证明中最锋利的工具。把``奇数''展开：存在整数 \(m,n\) 使得
\[
a=2m+1,\qquad b=2n+1.
\]
这步不是代数变形，这是把题目里模糊的``奇数''一词，翻译成了可以操作的等式。目标``偶数''也展开：存在某个整数 \(k\) 使得 \(a+b=2k\)。好，目标也清楚了——你只需要找出那个 \(k\)。

加一下：
\[
a+b=(2m+1)+(2n+1)=2(m+n+1).
\]
取 \(k=m+n+1\)，这是个整数。\(a+b=2k\)，符合偶数的定义。证完。

注意这段证明的骨架：假设给了你什么表示式（奇数定义），结论要求你产出什么表示式（偶数定义），中间几步代数只是把前者改写成后者的形状。直接证明的核心就是这个——把假设和结论都展开成定义，然后看看它们之间还差多远。多数时候，差的那段路就是几步基础代数。

\subsection{``任取''到底有多重要}

要证 \(\forall x\in D, P(x)\)，你开头选的那个 \(x\) 必须是真正任意的。你可以使用 \(x\) 属于论域 \(D\) 带来的任何性质，但不可以偷偷给它加一句``而且 \(x\) 还满足某某额外条件''。

比如证``任意实数的平方非负''。你可以分情况：\(x\ge0\) 和 \(x<0\)，两种情况合起来覆盖了所有实数，这没问题。但你不能一上来就写``令 \(x\) 为正实数''——这样一来，零和负数全被你排除在论证之外了，证明自然不完整。

一个证明的真实覆盖范围，从你写``任取 \(x\)''那句话的第一秒就定下来了。后面无论推导多漂亮，都不会变出一个比开头更大的论域。所以每次写``任取''，在心里默问一句：我真的允许读者塞进来任何一个符合条件的对象吗？

\subsection{集合相等：在两条方向间搬运元素}

现在换一道题：证明分配律
\[
A\cap(B\cup C)=(A\cap B)\cup(A\cap C).
\]

画个 Venn 图，三个圈交叉出来的区域确实一样。但图只能帮你发现这个等式的存在，不能构成证明——图上的点到底对应哪些元素、这些元素是否真的穿过每一层集合运算的门，只有逻辑语言能说清楚。

命题的形状是集合相等，那就走双向包含。

\textbf{方向一：左 \(\subseteq\) 右。}任取 \(x\in A\cap(B\cup C)\)。这顿饭拆开：\(x\in A\)，而且 \(x\in B\) 或者 \(x\in C\)。分叉了：
\begin{itemize}
\tightlist
\item 如果 \(x\in B\)，结合 \(x\in A\) 得 \(x\in A\cap B\)，那当然属于右边的并集；
\item 如果 \(x\in C\)，结合 \(x\in A\) 得 \(x\in A\cap C\)，同样属于右边的并集。
\end{itemize}
两条路都通向右边。左方向证完。

\textbf{方向二：右 \(\subseteq\) 左。}任取 \(x\) 在右边，它要么在 \(A\cap B\) 里，要么在 \(A\cap C\) 里。不管哪种情况，\(x\) 都在 \(A\) 里；同时在第一种情况里 \(x\in B\)，第二种里 \(x\in C\)，所以不管怎样 \(x\) 都在 \(B\cup C\) 里。于是 \(x\) 同时在 \(A\) 和 \(B\cup C\) 里，即 \(x\in A\cap(B\cup C)\)。

双向都通了。Venn 图帮你猜，元素追踪帮你证。两条腿走路，一条都别瘸。

\subsection{反例：一枚精确制导的子弹}

有些命题看着漂亮，实际是假的。推翻它们不需要长篇大论，一个反例就够。

看这个命题：``若 \(n^2\) 能被 4 整除，则 \(n\) 能被 4 整除。''拿 \(n=2\) 试试：\(2^2=4\)，4 当然能被 4 整除，但 2 本身不能被 4 整除。命题倒了。

不过要小心：\(n=3\) 并不能当反例。\(3^2=9\)，9 根本不能被 4 整除，前件就没成立。一个反例只有在满足了命题所有假设、却违反了结论的时候，才算有效的反驳。构造反例的最佳姿势是先把 \(P\land\neg Q\) 写出来——就是``假设全满足，结论偏不成立''——然后去找满足这个条件的对象。这能帮你过滤掉大量看起来可疑但实际不符合前件的候选。

反例的威力在于一击致命，但它的局限同样明确：反例只能推翻命题，反过来，举再多正面例子也证明不了全称命题。前一百万个整数都验证通过，也只是让猜想更可信而已，排除不了第一百万零一个出问题的可能。验证和证明之间，隔着一道全称量词的鸿沟。

\subsection{一段好证明应该让读者看到什么}

一个写得清楚的证明，不是把每一步初等代数都写在纸上——那是计算机干的事。但它必须让读者看清四件事：

\begin{enumerate}
\tightlist
\item
  讨论的对象是谁，论域是什么。一上来就让读者知道自己在哪个宇宙里。
\item
  你在用哪些假设。不要让读者一边读一边猜哪些条件被激活了。
\item
  从定义或定理得到了哪些关键的中间结论——这些是推理的骨架节点。
\item
  这些中间结论怎么拼在一起就正好够到结论。这是整段证明唯一的``啊哈''时刻，不能省。
\end{enumerate}

如果一句``显然''正好盖住了整段推理里唯一不显然的地方，那这个证明就失去了存在的意义。好的证明像一座桥：你不需要把每一颗铆钉都画出来，但桥墩和桥面必须清清楚楚地架在那里。

\subsection{形状匹配不是万能钥匙}

把命题按逻辑形状分类，是一个极其有效的起点。但它能给你的只是一个方向感——知道门朝哪边开，不等于知道钥匙长什么样。全称命题要证 \(P(x)\)，但 \(P(x)\) 内部可能嵌套着新的量词和新的逻辑结构；集合相等的两个方向，各自可能又是蕴含或析取。

形状告诉你``怎样才算证完了''，定义告诉你``手里有什么牌''。证明在这两者之间架桥。这个框架不会自动产生证明——它只是让你在迷宫里不至于连出口在哪都不知道。

延伸阅读：MIT Mathematics for Computer Science 的 Proofs 章节。
